paragraph{Background and Motivation.}

Laser ultrasound (LUS) has emerged as a promising technique for non-destructive testing (NDT) and evaluation, offering distinct advantages over conventional piezoelectric ultrasound methods. Unlike contact-based transducers, LUS enables non-contact generation and detection of ultrasonic waves through laser-induced thermoelastic effects~\cite{laser_us_ndt}. This non-contact nature eliminates coupling medium requirements, enabling inspection of rough, hot, or inaccessible surfaces where traditional methods are impractical. Furthermore, LUS systems produce broadband acoustic emissions with wide frequency content, facilitating high-resolution defect detection across diverse materials~\cite{laser_us_broadband}. The ability to perform rapid raster-scan or single-point measurements also makes LUS particularly attractive for industrial inspection pipelines where throughput is critical~\cite{laser_us_scanning}. These advantages have driven increasing adoption of LUS in aerospace, manufacturing, and infrastructure monitoring applications.

Despite its theoretical potential, practical deployment of LUS faces a fundamental challenge: the acquired signals exhibit severely degraded signal-to-noise ratio (SNR). Laser-generated ultrasound amplitudes are inherently limited by laser energy constraints and material absorption properties, while ambient electrical noise, laser source fluctuations, and environmental vibrations further contaminate the measured waveforms~\cite{low_snr_challenge}. This noise degradation is especially pronounced in field conditions or when inspecting低声学阻抗 materials, where the signal amplitudes may be orders of magnitude below noise floors. Consequently, reliable signal interpretation and defect characterization become difficult without effective denoising preprocessing.

Conventional digital signal processing (DSP) techniques have been extensively applied to LUS denoising, including Wiener filtering, wavelet thresholding, and block-matching 3D (BM3D) filtering~\cite{wiener_filter,wavelet_denoise,bm3d_denoise}. While these methods provide improvements under moderate noise conditions, they often fall short when SNR drops below practical thresholds. Wiener filters rely on accurate noise statistics that are difficult to estimate in real-world scenarios. Wavelet-based approaches, though effective at sparse representation, struggle to preserve fine signal features while suppressing non-Gaussian noise artifacts. BM3D, despite its strong performance on natural images, does not fully leverage the characteristic temporal structure of ultrasonic waveforms. These limitations motivate the development of data-driven learning-based approaches capable of adaptively capturing complex noise patterns and signal correlations inherent in LUS signals.
